
\sourcenote{Genre: admissible rule (tier 2). The family machinery is
stated in the \rulename{LiRPA} section's family template;
\rulename{ReLU-Pos} is this rule's mirror, and the shared boundary
argument is stated there.}

\begin{brecipe}[the \rulename{ReLU-Neg} instance]\label{rec:reluneg-instance}
The node, and its intended rewrite over the family's sign-split fact:
\[
\inferrule*[Right=\rulename{ReLU-Neg}]{
  \mathcal{T} \rhd \Phi \vdash h \leq 0 \\
  \ReLU_{\{h\}\to\{y\}} \in H
}{
  \Phi \vdash y = 0
}
\qquad\rightsquigarrow\qquad
\text{case analysis on }
(h \leq 0 \wedge y = 0) \vee (0 \leq h \wedge y = h).
\]
The \emph{inactive} branch discharges cleanly: its context carries the
conjunct $y = 0$, extracted by \rulename{Ax} and $\wedge$E --- the goal
verbatim. The \emph{boundary} branch is the mirror of
\rulename{ReLU-Pos}'s: its context carries $0 \leq h$ and $y = h$, the
premise supplies $h \leq 0$, and the goal $y = 0$ needs $h = 0$ --- the
antisymmetry fact that no catalogue derivation provides (the
enumeration argument is stated once, as
Remark~\ref{rem:relupos-antisym} in the \rulename{ReLU-Pos} section).
The rule is therefore admissible by reduction to verified machinery:
the shipped warrant is the direct lemma \texttt{sound\_reluNeg}, whose
boundary case is the same one-line antisymmetry step in the model.
\end{brecipe}

\begin{bprop}[semantic warrant]\label{prop:reluneg-sound}
Suppose the linkage holds and the premise sequent --- same premise
formula, conclusion the atom $h \leq 0$ --- is semantically valid.
Then the sequent concluding $y = 0$ is semantically valid.
\end{bprop}
\begin{proof}
Fix an admissible valuation $\sigma$ satisfying $\Phi$. The premise
eliminates entrywise: every entry $p$ of the preactivation tensor
satisfies $p \leq 0$. The linkage extraction
(\texttt{reluLinked\_facts}) gives, per entry, either $p \leq 0 \wedge
y = 0$ --- the goal entry directly --- or $0 \leq p \wedge y = p$, in
which case $p = 0$ by antisymmetry against the premise, so $y = p =
0$. The conclusion atom compares the output tensor entrywise against
the rank-$0$ constant $0$. That is \texttt{sound\_reluNeg}, step for
step.
\end{proof}

\begin{bremark}[what the premise is for, and what the rule feeds]
Without $\Phi \vdash h \leq 0$ the rule is false at any entry with
$h > 0$ (there $y = h \neq 0$). Downstream, the conclusion differs
from \rulename{ReLU-Pos}'s in kind: $y = 0$ pins the output to a
\emph{literal constant}, so \rulename{EqElim} can rewrite $y$ away
entirely --- dead-neuron elimination inside a certificate --- after
which ground arithmetic on the constant rides
\rulename{PConstExact}.
\end{bremark}

\begin{anchors}
Recipe \ref{rec:reluneg-instance} warrant & \texttt{sound\_reluNeg}\\
family linkage extraction & \texttt{reluLinked\_facts} (see the \rulename{LiRPA} section)\\
entrywise premise elimination & \texttt{QTensor.rel2\_const\_right\_elim}\\
the kink antisymmetry & the \texttt{le\_antisymm} step inside \texttt{sound\_reluNeg}\\
payload structure & \texttt{PReLUNegData} (premise proof, two wires, linkage)\\
\end{anchors}
